\documentclass[10pt,twocolumn]{article}
\usepackage[utf8]{inputenc}
\usepackage[T1]{fontenc}
\usepackage[a4paper,margin=0.55in]{geometry}
\usepackage{amsmath,amssymb}
\usepackage{enumitem}
\usepackage{hyperref}
\usepackage{setspace}
\setstretch{0.95}
\setlength{\parindent}{0pt}
\setlength{\parskip}{2.5pt}
\setlength{\columnsep}{14pt}
\setlength{\abovedisplayskip}{4pt}
\setlength{\belowdisplayskip}{4pt}
\setlength{\abovedisplayshortskip}{3pt}
\setlength{\belowdisplayshortskip}{3pt}
\setlist[itemize]{leftmargin=1.2em,itemsep=1pt,topsep=2pt}

\title{Reduced-Order Rayleigh-Taylor Dye Transport Model}
\author{}
\date{}

\begin{document}
\maketitle
\vspace{-2.0em}
\small

\section*{Purpose}
This document summarises the current \texttt{dye\_concentration\_evolution}
model used in the notebook. The model is not a full Navier-Stokes solver. It
is a reduced-order, state-dependent transport model in which the velocity field
evolves with the current dye and density structure.

\section*{Definitions}
The scalar field $\phi(x,z,t)\in[0,1]$ is the volume fraction of the low-density
dyed fluid. Therefore,
\begin{equation}
  \rho(\phi)=\phi\,\rho_{\text{light}}+(1-\phi)\,\rho_{\text{heavy}}.
\end{equation}
Here $x$ is horizontal and $z$ is vertical in the two-dimensional image plane.

\section*{Velocity Construction}
The model uses a prescribed background acceleration history $a(t)$ and a global
amplitude variable $f(t)$. The dye field first provides
the horizontal gradient
\begin{equation}
  \partial_x \phi \approx \frac{\phi_{i,j+1}-\phi_{i,j-1}}{2\Delta x}.
\end{equation}
This is used to construct a reduced-order baroclinic source term
\begin{equation}
  S_\psi = C_b\,f(t)\,\hat a(t)\,\partial_x \phi,
\end{equation}
where $C_b$ is the buoyancy-coupling coefficient and $\hat a(t)$ is the
normalised acceleration. The streamfunction is then obtained from the elliptic
closure
\begin{equation}
  \frac{\partial^2\psi}{\partial x^2}+\frac{\partial^2\psi}{\partial z^2}=S_\psi.
\end{equation}
Cell-centred velocities follow from
\begin{equation}
  u=\frac{\partial\psi}{\partial z}, \qquad
  w=-\frac{\partial\psi}{\partial x},
\end{equation}
with a vertical decay factor
\begin{equation}
  d(z)=\exp\left[-\left(\frac{z}{L_d}\right)^2\right].
\end{equation}
In practice, centred finite differences are used, so
\begin{equation}
  u_{\text{center}}=d(z)\,\partial_z\psi, \qquad
  w_{\text{center}}=-d(z)\,\partial_x\psi.
\end{equation}

\section*{Face Transport Velocity}
The centre velocities are interpolated to faces by arithmetic averaging:
\begin{equation}
  u_{\text{face}}\approx \mathrm{avg}(u_{\text{center}}), \qquad
  w_{\text{face}}\approx \mathrm{avg}(w_{\text{center}}).
\end{equation}
Two vertical corrections are then added. The buoyancy-drift term is
\begin{equation}
  w_{\text{drift}}
  =
  \texttt{rise\_velocity\_scale}\,g_w(z)\,(\phi-0.5),
\end{equation}
where $g_w(z)=0.25+0.75\left(1-(2z/L_z)^2\right)$. The unstable-jump sorting
term is
\begin{equation}
  \begin{aligned}
    w_{\text{sorting}}
    &=
    \texttt{sorting\_velocity\_scale}\,\texttt{flow\_strength}(t) \\
    &\qquad \cdot \max(\phi_{\text{below}}-\phi_{\text{above}},0).
  \end{aligned}
\end{equation}
The final vertical transport velocity is therefore
\begin{equation}
  w_{\text{face}}=
  w_{\text{baroclinic}}+w_{\text{drift}}+w_{\text{sorting}}.
\end{equation}

\section*{Dye Update}
The face velocities are converted into upwind face fluxes:
\begin{equation}
  F_x=\Delta t\,\Delta x\,(u_{\text{face}}\phi_{\text{upwind}}), \qquad
  F_z=\Delta t\,\Delta x\,(w_{\text{face}}\phi_{\text{upwind}}).
\end{equation}
For example,
\begin{equation}
  \phi_{\text{upwind},x,i,j+\frac{1}{2}}=
  \begin{cases}
    \phi_{i,j}, & u_{\text{face},i,j+\frac{1}{2}}\ge 0,\\
    \phi_{i,j+1}, & u_{\text{face},i,j+\frac{1}{2}}<0.
  \end{cases}
\end{equation}
The dye field is then updated by conservative advection plus diffusion:
\begin{equation}
  \phi^\ast=
  \phi^n+
  \frac{F_{x,\text{left}}+F_{z,\text{bottom}}-F_{x,\text{right}}-F_{z,\text{top}}}
       {\Delta x\Delta z}
  +\kappa_{\mathrm{eff}}\Delta t\nabla^2\phi.
\end{equation}
After this explicit step, the code clips $\phi$ into $[0,1]$, applies weak
mass-preserving sharpening, and enforces the global invariants.

\section*{Interpretation}
The model therefore follows the chain
\[
\phi^n \rightarrow S_\psi \rightarrow \psi \rightarrow (u,w)_{\text{center}}
\rightarrow (u,w)_{\text{face}} \rightarrow (F_x,F_z) \rightarrow \phi^{n+1}.
\]
It is more physical than a fixed placeholder sine-wave velocity field because
the transport field depends on the current density structure, but it remains a
reduced-order RT-like closure rather than a full CFD solution.

\end{document}
